\subsection{Linked Census Microdata}\label{sec:data-IPUMS}

I collect economic and demographic data on the populations of interest from the full count decenial censuses taken in 1940 and 1950
\citep{rugglesIPUMSAncestryFull2024}.
These censuses asked about the wages and salary income earned by respondents
which I use as the main labor market outcome in my analysis of the impact of internment and forced removal on internees.
The censuses also provide information on the employment status of respondents;
demographic variables such as age, sex, race, etc;
and the county of residence when the census survey was administered.


\subsection{War Relcation Authority Records}

The WRA kept detailed records of the process of relocating and incarcerating the West Coast Japanese population.
I use the dataset from WRA Form 26: Evacuee Summary Data in the collection of Records About Japanese Americans Interned During World War II in the National Archives.
The original paper form from which this digitized version is coded included each internee's name, previous address in 1942, the assembly center and internment camp they were assigned to,
their educational attainment, occupation, and demographics including sex, birth year, and birthplace.

Table \ref{tbl:bal-stat} reports a summary of the main demographic variables of internees compared to the full 1940 Japanese population.
Internees represented a large majority (86\%) of the Japanese population in the US in 1940.
Internees were representative of the population in terms of their education, gender composition, and country of birth.
The 65\% of internees born in the US represent the generation known as the ``Nisei'', or second generation after their parents, 
the ``Issei'', who make up most of the remaining 35\% of internees.
The average Issei was born in 1890, and would have been in their early 50s during internment,
while the average Nisei was born in 1927 and would have been teenagers when they were interned.

Figure \ref{fig:internee-ages} shows the distribution of the age of internees at the time of the evacutation order in 1942.
This figure reveals that there was a large variation in the ages of internees and that the ten camps were fairly balanced on their age composition.
The most prominant age range in this figure is made up of young adults between the ages of 15-25. 

The WRA records also provide a matching between the county of residence of all internees to the first camp they were placed.
Table \ref{tbl:cmp-cty} groups all internees into categories of their previous address to their assigned camp.
The largest 20 counties with more than 1000 interned residents are shown broken down by each of the ten camps.

Among the main internment counties, the major West Coast urban areas of Los Angeles, Seattle, Sacramento, San Franciso, and Portland.
However, there are also a large number of internees from more rural counties in California including Fresno, San Joaquin, Placer, and Tulare.
Interned Japanese came from a variety of backgrounds shown here by the range of urban and rural counties,
and also reflected in the range of occupations held by internees.
